\section{Matching across cuffs: positions and directions}
\label{sec:matching}

The pants viewpoint suggests that orthogeodesic pieces should glue when
their boundary data match. There are two levels of matching.

First consider a positive-length cuff $a>0$, with $m\geq1$ distinct
endpoint positions $x_1,\ldots,x_m$ and $y_1,\ldots,y_m$ on the two sides.
Choose compatible
coordinates in $\RR/a\ZZ$, reversing one induced boundary orientation
as needed. For a prescribed cyclic matching $\sigma$, a twist $\tau$
must satisfy
\begin{equation}
 y_{\sigma(j)}-x_j\equiv\tau\pmod a
       \qquad(1\le j\le m).
 \label{eq:endpointmatch}
\end{equation}
One endpoint determines $\tau$ modulo $a$; the remaining $m-1$ equations
require the cyclic spacings to agree. Integer winding data must also be
retained if the curve is marked, rather than only considered modulo a
twist. These congruences are necessary and sufficient for the prescribed
position matching. They do not supply the direction data of the pieces.

For example, let the cuff length be $6$ and take positions
$x=(1,3)$ and $y=(2,4)$ in the chosen cyclic order. Both
differences are $1$, so twist $\tau=1$ matches them. Replacing
$y$ by $(2,5)$ gives differences $1$ and $2$, so no single
twist realizes that prescribed matching. Matching one endpoint
has not determined the rest of the chain.

\begin{figure}[htbp]
\centering
\begin{tikzpicture}[font=\small,>=Latex]
 \draw[mpgray] (0,1)--(6,1);\draw[mpgray] (0,-.4)--(6,-.4);
 \foreach \x in {0,6}{\draw (\x,.9)--(\x,1.1);\draw (\x,-.5)--(\x,-.3);}
 \node[left] at (0,1) {first side};\node[left] at (0,-.4) {second side};
 \foreach \x/\lab in {1/{x_1=1},3/{x_2=3}}
   {\fill[mpblue] (\x,1) circle(2pt);\node[above] at (\x,1.1) {$\lab$};}
 \foreach \x/\lab in {2/{y_1=2},4/{y_2=4}}
   {\fill[mpteal] (\x,-.4) circle(2pt);\node[below] at (\x,-.5) {$\lab$};}
 \draw[->,mpochre] (1,.85)--(2,-.25);
 \draw[->,mpochre] (3,.85)--(4,-.25);
 \node[right] at (6,.3) {$\tau=1\pmod6$};
\end{tikzpicture}
\caption*{Position matching on a cuff of length $6$, drawn as two
intervals whose endpoints $0$ and $6$ are identified on each side.
The arrows show the same translation for both marked endpoints.
They indicate the gluing correspondence, not geodesic directions;
tangent matching must still be checked.}
\end{figure}

Second, orient a chain by its direction of travel and measure all
tangential components in the common cuff coordinate just chosen.
A transverse unit-speed geodesic segment has a signed tangential velocity
component $p\in(-1,1)$ at an endpoint. If $S$ is the sum of segment lengths, the first
variation at a glued endpoint has the form
\begin{equation}
 \frac{\partial S}{\partial x}=p_{\rm in}-p_{\rm out}.
 \label{eq:momentummatch}
\end{equation}
Here $p_{\rm in}$ is the terminal tangential component of the incoming
segment and $p_{\rm out}$ the initial component of the outgoing segment.
The pieces must approach and leave on the opposite sides of the cut.
Equality of their tangential components then also matches the normal
components, whose magnitudes are $\sqrt{1-p^2}$ and whose signs are fixed
by this crossing condition. Stationarity therefore imposes full tangent
matching as well as position matching. For
orthogonal segments the tangential components vanish, so smooth gluing
is automatic once the endpoints match. A generic closed geodesic does
not satisfy this orthogonality condition along an arbitrarily chosen
pants decomposition.

At one of those matched positions, incoming and outgoing tangential
components $p_{\rm in}=1/2$ and $p_{\rm out}=-1/2$ give
$\partial S/\partial x=1$, so the chain has a corner and is
not stationary. If both are $1/2$ and the pieces cross to opposite
sides, their normal components have magnitude $\sqrt3/2$ and
fit as well. Position and direction are separate tests.

Equivalently, boundary data live in a phase space $(s,p)$: endpoint
position and tangential component. The length of an arc is a generating
function, with
\[
 \dd L=-p_s\,\dd s+p_t\,\dd t,
 \qquad\text{and cuff gluing acts by }(s,p)\longmapsto(s+\tau,p).
\]
A nonperipheral periodic chain with these position and direction matches
is smooth and locally geodesic, hence is its closed geodesic
representative. Conversely, cutting a transverse closed geodesic gives
such a chain. Tangent segments lying on a cuff require separate treatment
and are not included in the transverse phase space.
The orthospectrum records the special $p=0$ representatives. It is
excellent input for gap identities and local geometry, but endpoint
positions and directions are needed to reconstruct arbitrary glued curves.

\paragraph{Computational support.}
Group G8 of \path{experiments/quantum_geodesics/README.md}
checks finite endpoint/winding matches and a directional stationarity
example. These distinguish position matching from the tangential
momentum condition; they do not supply a general geometric gluing
solver. The result path is \path{artifacts/quantum_geodesics/result.json},
with claim pointers in \path{book/claims/sections17-22.json}.
